\chapter{Introduzione al Web}

\section{L'architettura di una web app}

Prima di entrare nei dettagli conviene avere in testa la mappa del territorio.
Una \term{web application} moderna è quasi sempre un sistema distribuito
composto da due metà che girano su macchine diverse e comunicano attraverso la
rete.

\begin{itemize}
  \item Il \term{front-end} è la parte eseguita nel \textbf{browser} dell'utente.
        È ciò che l'utente vede e con cui interagisce: la struttura della pagina,
        il suo aspetto, il suo comportamento interattivo.
  \item Il \term{back-end} è la parte eseguita su uno o più \textbf{server}. Non
        è visibile all'utente, ma è ciò che custodisce i dati, applica le regole
        di business, autentica gli utenti e produce le risposte.
\end{itemize}

Le due metà si parlano quasi sempre tramite \term{HTTP}. Chi sviluppa entrambe
le parti viene chiamato \textit{full-stack developer}.

\begin{figure}[H]
  \centering
  \begin{tikzpicture}[font=\Lato]
    \node[wtnode, minimum width=3.6cm, minimum height=1.7cm] (br)
      {\textbf{Web Browser}\\[3pt]\scriptsize\textsc{front-end}};
    \node[wtnode, minimum width=3.6cm, minimum height=1.7cm, right=4.8cm of br] (sv)
      {\textbf{Web Server}\\[3pt]\scriptsize\textsc{back-end}};
    \node[wtnodeplain, minimum width=2.2cm, minimum height=1.7cm, right=1.2cm of sv] (db)
      {Database};

    \draw[wtarrow]    ([yshift=7pt]br.east)  -- node[wtlab, above]{richiesta HTTP} ([yshift=7pt]sv.west);
    \draw[wtarrowacc] ([yshift=-7pt]sv.west) -- node[wtlab, below]{risposta HTTP}  ([yshift=-7pt]br.east);
    \draw[wtarrow, {Stealth[length=5pt]}-{Stealth[length=5pt]}] (sv.east) -- (db.west);

    \node[wtnodeplain, below=0.7cm of br, text width=3.0cm]
      {HTML, CSS,\\ JavaScript,\\ Angular / React};
    \node[wtnodeplain, below=0.7cm of sv, text width=3.0cm]
      {Node.js, Express,\\ PHP, Java,\\ Python};
  \end{tikzpicture}
  \caption{Architettura di massima di una web application full-stack.}
  \label{fig:fullstack}
\end{figure}

\info{Anche le app mobili sono \guillemotleft web\guillemotright}{
  Quando si usa l'app di una banca, di un social network o di un servizio di
  consegne, quasi sempre quell'app non fa altro che inviare richieste HTTP a un
  server remoto e mostrare le risposte. Molto spesso, inoltre, la sua interfaccia
  grafica è realizzata con le stesse tecnologie delle pagine web (HTML, CSS e
  JavaScript) racchiuse in un contenitore nativo. Imparare le tecnologie web,
  quindi, serve ben oltre il browser.}

\section{Internet e il World Wide Web}

È molto comune usare le parole \guillemotleft Internet\guillemotright\ e \guillemotleft Web\guillemotright\ come se fossero
sinonimi. Non lo sono.

\dfn{Internet}{
  \term{Internet} è una rete globale di calcolatori interconnessi che si
  scambiano informazioni utilizzando i \textit{protocolli Internet} (la famiglia
  TCP/IP). Le sue origini risalgono ad \textbf{ARPANET} (1969).}

\dfn{World Wide Web}{
  Il \term{World Wide Web} (WWW, o semplicemente \term{Web}) è un
  \textbf{sottoinsieme} di Internet: è un sistema di documenti ipertestuali
  interconnessi, collegati fra loro tramite collegamenti ipertestuali
  (\textit{hyperlink}). È stato inventato da Sir \textbf{Tim Berners-Lee} nei
  primi anni Novanta. I suoi due componenti fondamentali sono il protocollo
  \textbf{HTTP} e il linguaggio \textbf{HTML}.}

La differenza è sostanziale: Internet è l'\textit{infrastruttura} (cavi, router,
indirizzi IP, protocolli di trasporto); il Web è uno dei tanti \textit{servizi}
che vi girano sopra. La posta elettronica (SMTP, IMAP), il trasferimento file
(FTP), le videochiamate, i giochi online e i sistemi di messaggistica usano
tutti Internet, ma non sono il Web.

\begin{figure}[H]
  \centering
  \begin{tikzpicture}[font=\Lato\small]
    \draw[line width=1pt, draw=primary!60!black, fill=primary!6!white]
      (0,0) ellipse (4.8cm and 2.3cm);
    \node[font=\Lato\bfseries, text=primary!70!black] at (0,1.75) {INTERNET};

    \draw[line width=1pt, draw=accent!60!black, fill=accent!10!white]
      (-2.1,-0.2) ellipse (2.2cm and 1.25cm);
    \node[font=\Lato\bfseries\footnotesize, text=accent!70!black] at (-2.1,0.45) {WEB};
    \node[font=\Lato\scriptsize, align=center, text=accent!45!black] at (-2.1,-0.5)
      {HTTP\\ HTML};

    \node[font=\Lato\scriptsize, align=center, text=primary!40!black] at (2.3,0.1)
      {e-mail (SMTP, IMAP)\\ FTP \quad DNS\\ SSH \quad VoIP};
  \end{tikzpicture}
  \caption{Il Web è uno dei servizi che viaggiano su Internet, non Internet stesso.}
  \label{fig:internet-web}
\end{figure}

\section{Documenti ipertestuali}

I documenti a cui siamo abituati (un file di testo, una lettera) sono
essenzialmente \textbf{sequenze di caratteri}: si leggono dall'inizio alla fine,
in modo lineare.

\dfn{Ipertesto}{
  Un \term{ipertesto} è un documento che contiene anche \term{collegamenti}
  (\textit{hyperlink}) verso altri contenuti: altri ipertesti, documenti,
  immagini o altri file multimediali.}

È proprio questa la rivoluzione del Web: la lettura smette di essere lineare. Da
un documento si può saltare a un altro documento, magari ospitato su un server
dall'altra parte del mondo, semplicemente attivando un collegamento.

\begin{esempio}{Un ipertesto minimale}
Supponiamo di leggere questo frammento tratto dal \textit{Libro della
Programmazione}:

\begin{browser}[bookofprogramming.com/wisdom.html]
  {\Lora\bfseries\large Frammento dal Libro della Programmazione}\\[6pt]
  Uno studente chiese: \guillemotleft I programmatori di un tempo usavano solo macchine
  semplici e nessun linguaggio di programmazione, eppure scrivevano programmi
  bellissimi. Perché noi usiamo macchine complicate e linguaggi di
  programmazione?\guillemotright\ \textcolor{blue}{\underline{Fu-Tzu}} rispose: \guillemotleft I costruttori
  di un tempo usavano solo bastoni e argilla, eppure costruivano capanne
  bellissime\guillemotright.
\end{browser}

La parola \textit{Fu-Tzu} è sottolineata e colorata: è un collegamento. Chi la
attiva viene portato a un altro documento, \code{fu-tzu.html}, che descrive il
personaggio. I due documenti possono trovarsi su server diversi, e nessuno dei
due contiene il testo dell'altro: contiene solo il suo \textit{indirizzo}.
\end{esempio}

\section{HTTP: il protocollo del Web}

\dfn{HTTP}{
  \term{HTTP} (\textit{HyperText Transfer Protocol}) è un protocollo di
  \textbf{livello applicativo}, costruito sopra TCP/IP, che costituisce la
  spina dorsale del World Wide Web. Nato per trasferire ipertesti, oggi è usato
  per trasferire risorse di ogni tipo (immagini, video, dati JSON, file binari).}

Il funzionamento di HTTP è basato sul modello \term{client-server} ed è
sorprendentemente semplice:

\begin{enumerate}
  \item il \textbf{client} (tipicamente un browser) invia una \term{richiesta}
        per una particolare risorsa;
  \item il \textbf{server} elabora la richiesta e invia una \term{risposta}.
\end{enumerate}

Le risorse sono identificate dalla loro \term{URL} (\textit{Uniform Resource
Locator}).

\begin{figure}[H]
  \centering
  \begin{tikzpicture}[font=\Lato]
    \node[wtnode, minimum width=2.8cm, minimum height=1.3cm] (cl) {\textbf{Client}};
    \node[wtnode, minimum width=2.8cm, minimum height=1.3cm, right=5.8cm of cl] (sr) {\textbf{Server}};
    \draw[wtarrow]    ([yshift=8pt]cl.east)  -- node[wtlab, above]{\textbf{Richiesta} HTTP} ([yshift=8pt]sr.west);
    \draw[wtarrowacc] ([yshift=-8pt]sr.west) -- node[wtlab, below]{\textbf{Risposta} HTTP}  ([yshift=-8pt]cl.east);
  \end{tikzpicture}
  \caption{Il ciclo richiesta--risposta, unità elementare di ogni interazione HTTP.}
  \label{fig:req-res}
\end{figure}

\info{Il client parla per primo, sempre}{
  In HTTP l'iniziativa è \textbf{sempre} del client: il server non può
  contattare spontaneamente il browser per dirgli che qualcosa è cambiato.
  Questo vincolo, apparentemente banale, è all'origine di molte soluzioni
  tecniche che vedremo più avanti (\textit{polling}, \textit{long polling},
  WebSocket, \textit{Server-Sent Events}): tutte servono a simulare una
  comunicazione avviata dal server, su un protocollo che non la prevede.}

\section{URL: individuare una risorsa}

Ogni risorsa sul Web è identificata da una URL. Una URL non è una stringa
qualunque: ha una struttura precisa, e ogni sua parte ha un ruolo.

\begin{figure}[H]
  \centering
  \begin{tikzpicture}[
      part/.style={font=\FiraCode\large, inner sep=1pt, outer sep=0pt},
      lab/.style={font=\Lato\scriptsize, align=center, inner sep=2pt}]
    \node[part, text=codekey]                       (sch)  {https://};
    \node[part, text=codetag,    right=0pt of sch]  (dom)  {www.informatica.it};
    \node[part, text=codenum,    right=0pt of dom]  (port) {:4242};
    \node[part, text=codestring, right=0pt of port] (path) {/corsi/tecweb.html};

    \node[lab, text=codekey, above=18pt of sch]  (lsch)  {Schema\\(protocollo)};
    \node[lab, text=codenum, above=18pt of port] (lport) {Porta};
    \draw[wtdash] (lsch.south)  -- (sch.north);
    \draw[wtdash] (lport.south) -- (port.north);

    \node[lab, text=codetag,    below=18pt of dom]  (ldom)  {Nome di dominio};
    \node[lab, text=codestring, below=18pt of path] (lpath) {Percorso};
    \draw[wtdash] (ldom.north)  -- (dom.south);
    \draw[wtdash] (lpath.north) -- (path.south);
  \end{tikzpicture}
  \caption{Anatomia di una URL.}
  \label{fig:url}
\end{figure}

\subsection{Schema}

Lo \term{schema} specifica il protocollo da usare per accedere alla risorsa. I
due più comuni sul Web sono \code{http} e \code{https}.

\term{HTTPS} (\textit{HTTP Secure}) è semplicemente HTTP trasportato su una
connessione \textbf{cifrata}. La cifratura è ottenuta tramite \term{TLS}
(\textit{Transport Layer Security}) e gioca un ruolo importante nel mitigare
alcune classi di attacchi alle applicazioni web, che vedremo nella lezione
dedicata alla sicurezza.

\warning{HTTPS non è un optional}{
  Senza TLS, chiunque si trovi lungo il percorso di rete (l'amministratore della
  rete Wi-Fi del bar, il provider, un attaccante sulla stessa rete) può leggere
  \textbf{in chiaro} tutto ciò che viaggia: password, cookie di sessione, dati
  personali. Oggi tutti i browser marcano come \guillemotleft non sicuri\guillemotright\ i siti
  serviti in HTTP semplice, e diverse funzionalità del browser
  (geolocalizzazione, fotocamera, \textit{service worker}) sono disponibili solo
  su HTTPS.}

\subsection{Nome di dominio}

Il \term{nome di dominio} identifica il server web che ospita la risorsa. È
composto da più parti separate da punti e va letto \textbf{da destra verso
sinistra}:

\begin{itemize}
  \item la prima parte da destra (\code{.it}) è il \term{Top Level Domain} (TLD).
        La \textit{Internet Assigned Numbers Authority} (IANA) mantiene l'elenco
        ufficiale dei TLD esistenti;
  \item la seconda (\code{informatica}) è il \term{Secondary Level Domain} (SLD);
  \item le parti ulteriori definiscono i \term{sottodomini}, usati per
        distinguere contenuti diversi ospitati sullo stesso dominio.
\end{itemize}

\begin{esempio}{Lettura di alcuni domini}
\begin{center}
\renewcommand{\arraystretch}{1.3}
\begin{tabular}{@{}lll@{}}
  \textbf{Dominio} & \textbf{TLD} & \textbf{Come si legge} \\
  \code{blog.mozilla.org}           & \code{.org} & sottodominio \code{blog} di \code{mozilla} \\
  \code{informatica.dieti.unina.it} & \code{.it}  & \code{informatica} dentro \code{dieti} dentro \code{unina} \\
  \code{luistar.github.io}          & \code{.io}  & sottodominio \code{luistar} di \code{github} \\
\end{tabular}
\end{center}
\end{esempio}

\subsection{Porta}

La \term{porta} indica su quale porta TCP contattare il server.

\begin{itemize}
  \item Può essere \textbf{omessa} se il server usa le porte standard.
  \item La porta standard per HTTP è la \code{80}, per HTTPS la \code{443}.
  \item Deve essere specificata esplicitamente se il server usa una porta non
        standard (nell'esempio, la \code{4242}).
\end{itemize}

\info{Perché in sviluppo si vede sempre \code{localhost:3000}}{
  Sulle porte inferiori a 1024 (fra cui la 80 e la 443) i sistemi Unix
  richiedono privilegi di amministratore per mettersi in ascolto. Per questo i
  server di sviluppo usano quasi sempre porte alte come 3000, 8000 o 8080, e
  l'indirizzo va scritto per esteso: \code{http://localhost:3000}.}

\subsection{Percorso}

Il \term{percorso} (\textit{path}) indica la posizione specifica, sul server, in
cui si trova la risorsa.

Il percorso è tipicamente \textbf{relativo a una directory radice}
(\textit{web root} o \textit{document root}) configurata sul server: il server
serve solo i file contenuti in quella directory, e nulla al di fuori di essa. La
ragione è evidente: non vogliamo che \textbf{tutti} i file della macchina siano
accessibili dal Web.

\error{Path traversal}{
  Un server configurato male, o un'applicazione che costruisce percorsi
  concatenando input dell'utente, può essere indotto a servire file fuori dalla
  document root con richieste del tipo \code{/../../../etc/passwd}. Questo
  attacco si chiama \textit{path traversal} ed è uno dei motivi per cui non
  bisogna \textbf{mai} fidarsi dell'input dell'utente. Ne riparleremo nel
  capitolo dedicato alla sicurezza.}

Le URL possono contenere anche \term{parametri di query} e \term{ancore}: li
vedremo nel prossimo capitolo, insieme ai form HTML.

\section{I messaggi HTTP}

Un'interazione HTTP è fatta di due tipi di messaggi: la \term{richiesta}
(\textit{request}) e la \term{risposta} (\textit{response}). Entrambi hanno la
stessa struttura a quattro parti.

\begin{figure}[H]
  \centering
  \begin{tikzpicture}[font=\Lato\footnotesize,
      strip/.style={draw=primary!60!black, line width=0.7pt, minimum width=5.4cm,
                    minimum height=0.75cm, align=center, anchor=north},
      striph/.style={strip, fill=primary!22!white, font=\Lato\footnotesize\bfseries},
      stripb/.style={strip, fill=primary!7!white},
      stripe/.style={strip, fill=neutral!10!white, minimum height=0.42cm,
                     font=\Lato\scriptsize, text=neutral!30!black}]

    \node[font=\Lato\bfseries, text=primary!75!black] at (0,0.95) {RICHIESTA};
    \node[striph] (r1) at (0,0.5) {Verbo \quad URI \quad Versione HTTP};
    \node[stripb, below=0pt of r1] (r2) {Header della richiesta};
    \node[stripe, below=0pt of r2] (r3) {riga vuota};
    \node[stripb, below=0pt of r3, minimum height=1.1cm] (r4)
      {Corpo della richiesta\\ \scriptsize(opzionale)};

    \begin{scope}[xshift=6.4cm]
      \node[font=\Lato\bfseries, text=accent!75!black] at (0,0.95) {RISPOSTA};
      \node[striph, draw=accent!60!black, fill=accent!20!white] (s1) at (0,0.5)
        {Versione HTTP \quad Stato};
      \node[stripb, draw=accent!60!black, fill=accent!7!white, below=0pt of s1] (s2)
        {Header della risposta};
      \node[stripe, draw=accent!60!black, below=0pt of s2] (s3) {riga vuota};
      \node[stripb, draw=accent!60!black, fill=accent!7!white, below=0pt of s3,
            minimum height=1.1cm] (s4) {Corpo della risposta\\ \scriptsize(opzionale)};
    \end{scope}
  \end{tikzpicture}
  \caption{Struttura dei due tipi di messaggio HTTP.}
  \label{fig:http-msg}
\end{figure}

\subsection{Metodi (o verbi) di richiesta}

Il \term{metodo} indica l'azione che si desidera compiere sulla risorsa
indicata. I metodi più comuni sono:

\begin{center}
\renewcommand{\arraystretch}{1.4}
\begin{tabular}{@{}l@{\hspace{1.1cm}}>{\raggedright\arraybackslash}p{10cm}@{}}
  \textbf{Metodo} & \textbf{Descrizione} \\
  \code{GET}    & Recupera (una rappresentazione di) una risorsa. \\
  \code{POST}   & Invia nuovi dati alla risorsa specificata. \\
  \code{PUT}    & Sostituisce la risorsa corrente con il contenuto inviato. \\
  \code{DELETE} & Cancella la risorsa specificata. \\
\end{tabular}
\end{center}

L'elenco completo dei metodi HTTP è disponibile sulla documentazione MDN. Ne
esistono altri (\code{HEAD}, \code{OPTIONS}, \code{PATCH}, \code{TRACE},
\code{CONNECT}) che incontreremo parlando di API REST.

\info{Il verbo è una promessa, non un obbligo}{
  Nulla nel protocollo impedisce di scrivere un server in cui una \code{GET}
  cancella dei dati. È però una pessima idea: browser, \textit{proxy} e motori
  di ricerca assumono che \code{GET} sia \term{sicura} (non modifica nulla) e
  \term{idempotente} (ripeterla non cambia il risultato), e sono liberi di
  ripeterla o di metterla in cache. Rispettare la semantica dei verbi non è
  formalismo: è ciò che rende prevedibile l'intera infrastruttura del Web.}

\subsection{Header}

Gli \term{header} (intestazioni) sono il modo con cui si trasmettono
informazioni aggiuntive, sia nelle richieste sia nelle risposte.

Un header è formato dal suo \textbf{nome} (non sensibile a maiuscole e
minuscole), seguito dai due punti e dal suo \textbf{valore}:

\begin{center}
  \code{NOME\_HEADER: valore}
\end{center}

La IANA mantiene una lista degli header permanenti e provvisori; è inoltre
possibile definire header \textbf{personalizzati}. Un riferimento completo e
aggiornato si trova nella documentazione MDN.

\subsection{Esempio di richiesta}

\begin{esempio}{Una richiesta HTTP completa}
\begin{lstlisting}[style=http, name=Richiesta HTTP]
GET /wisdom/grain.txt HTTP/1.1
Host: bookofprogramming.com
User-Agent: Mozilla/5.0
Accept: text/plain
Accept-Language: en-us
Connection: keep-alive
\end{lstlisting}

Leggiamola riga per riga:

\begin{itemize}
  \item \code{GET} è il \textbf{metodo}: si sta chiedendo di leggere una risorsa;
  \item \code{/wisdom/grain.txt} è l'\textbf{URI}, cioè il percorso della risorsa
        sul server;
  \item \code{HTTP/1.1} è la \textbf{versione} del protocollo;
  \item \code{Host} indica \textit{quale} sito si sta chiedendo. È l'unico header
        obbligatorio in HTTP/1.1, e serve perché un solo server può ospitare
        molti siti diversi (lo vedremo nel prossimo capitolo);
  \item \code{User-Agent} identifica il programma che effettua la richiesta;
  \item \code{Accept} e \code{Accept-Language} dichiarano quali formati e quali
        lingue il client preferisce ricevere;
  \item \code{Connection: keep-alive} chiede di non chiudere la connessione TCP
        dopo la risposta, così da riutilizzarla per le richieste successive;
  \item segue una \textbf{riga vuota}, che segnala la fine degli header;
  \item il \textbf{corpo} è assente: una \code{GET} normalmente non ne ha uno.
\end{itemize}
\end{esempio}

\subsection{Codici di stato della risposta}

Il \term{codice di stato} (\textit{status code}) indica se la richiesta è stata
portata a termine con successo. I codici sono raggruppati in cinque classi,
identificate dalla prima cifra.

\begin{figure}[H]
  \centering
  \begin{tikzpicture}[font=\Lato\footnotesize,
      cls/.style={rounded corners=3pt, line width=0.8pt, minimum width=2.6cm,
                  minimum height=1.0cm, align=center, anchor=north,
                  font=\Lato\scriptsize\bfseries},
      det/.style={align=center, font=\FiraCode\scriptsize, anchor=north,
                  text=neutral!25!black}]
    \node[cls, draw=neutral!55!white,  fill=neutral!12!white]         (c1) at (0.00,0) {Informational\\ 100--199};
    \node[cls, draw=green-gradient-6,  fill=green-gradient-1!55!white] (c2) at (2.75,0) {Success\\ 200--299};
    \node[cls, draw=azure-gradient-6,  fill=azure-gradient-1!55!white] (c3) at (5.50,0) {Redirection\\ 300--399};
    \node[cls, draw=yellow-gradient-6, fill=yellow-gradient-1!55!white](c4) at (8.25,0) {Client Error\\ 400--499};
    \node[cls, draw=red-gradient-6,    fill=red-gradient-1!45!white]   (c5) at (11.00,0) {Server Error\\ 500--599};

    \node[det, below=6pt of c1] {100 Continue};
    \node[det, below=6pt of c2] {200 OK\\ 201 Created};
    \node[det, below=6pt of c3] {301 Moved\\ 304 Not Modified};
    \node[det, below=6pt of c4] {400 Bad Request\\ 403 Forbidden\\ 404 Not Found};
    \node[det, below=6pt of c5] {500 Internal Error\\ 503 Unavailable};
  \end{tikzpicture}
  \caption{Le cinque classi di codici di stato HTTP.}
  \label{fig:status}
\end{figure}

\warning{4xx contro 5xx}{
  La distinzione fra le due classi di errore è importante e viene spesso
  sbagliata. Un \code{4xx} dice: \guillemotleft la colpa è della tua richiesta\guillemotright\
  (risorsa inesistente, permessi mancanti, dati malformati), e ripeterla
  identica non servirà a nulla. Un \code{5xx} dice: \guillemotleft la richiesta era
  plausibile, ma io ho un problema\guillemotright, e riprovare più tardi ha senso.}

\subsection{Esempio di risposta}

\begin{esempio}{Una risposta HTTP completa}
\begin{lstlisting}[style=http, name=Risposta HTTP]
HTTP/1.1 200 OK
Content-Type: text/plain
Content-Length: 153

Fu-Tzu disse: quando tagli il legno contro venatura, serve molta
forza. Quando programmi contro la venatura di un problema, serve
molto codice.
\end{lstlisting}

\begin{itemize}
  \item \code{HTTP/1.1} è la versione del protocollo;
  \item \code{200} è il codice di stato e \code{OK} il suo messaggio testuale;
  \item \code{Content-Type} dichiara il \textbf{tipo} di contenuto restituito
        (qui testo semplice; per una pagina web sarebbe \code{text/html});
  \item \code{Content-Length} ne dichiara la \textbf{lunghezza} in byte, così che
        il client sappia quando il corpo è finito;
  \item dopo la riga vuota comincia il \textbf{corpo} della risposta, cioè il
        contenuto vero e proprio della risorsa.
\end{itemize}
\end{esempio}

\section{HTTP è un protocollo stateless}

\dfn{Statelessness}{
  HTTP è un protocollo \term{stateless} (\textit{senza stato}): ogni richiesta è
  \textbf{indipendente} da quelle che l'hanno preceduta, e il server non conserva
  alcuna informazione sulle richieste precedenti di un dato client.}

Questa è una delle decisioni di progetto più importanti (e più controintuitive)
del Web. Il vantaggio è enorme in termini di \textbf{scalabilità}: se il server
non deve ricordarsi nulla, allora due richieste dello stesso utente possono
essere servite da due macchine diverse, e aggiungere macchine è sufficiente per
servire più utenti.

Lo svantaggio è altrettanto evidente: senza stato non esiste il concetto di
\guillemotleft utente collegato\guillemotright. Se effettuiamo il login e poi chiediamo un'altra
pagina, per il server siamo dei perfetti sconosciuti.

Per questo servono meccanismi \textbf{costruiti sopra} il protocollo per simulare
una comunicazione con stato. Il più noto è quello dei \term{cookie}, che
studieremo fra qualche lezione.

\section{Cosa succede quando si apre una pagina}

Mettiamo insieme i pezzi. Quando digitiamo
\code{http://example.com/dir/file.txt} nella barra degli indirizzi del browser,
accade la sequenza seguente.

\begin{figure}[H]
  \centering
  \begin{tikzpicture}[font=\Lato\small,
      hdr/.style={wtnode, minimum width=2.9cm, minimum height=0.9cm},
      life/.style={draw=neutral!45!white, line width=0.7pt, dashed},
      msg/.style={-{Stealth[length=6pt,width=5pt]}, line width=0.9pt,
                  draw=primary!70!black},
      rsp/.style={-{Stealth[length=6pt,width=5pt]}, line width=0.9pt,
                  draw=accent!70!black},
      ml/.style={font=\Lato\scriptsize, midway, above, inner sep=2.5pt,
                 fill=white, text=neutral!15!black},
      num/.style={circle, fill=accent!75!black, text=white, inner sep=1.5pt,
                  font=\Lato\bfseries\tiny}]

    \def\xb{0} \def\xd{5.3} \def\xs{10.6}

    \node[hdr] (B) at (\xb,0) {\textbf{Browser}};
    \node[hdr, draw=orange-gradient-6, fill=orange-gradient-1!45!white]
      (D) at (\xd,0) {\textbf{Server DNS}};
    \node[hdr] (S) at (\xs,0) {\textbf{Server HTTP}};

    \draw[life] (B.south) -- (\xb,-5.6);
    \draw[life] (D.south) -- (\xd,-5.6);
    \draw[life] (S.south) -- (\xs,-5.6);

    \draw[msg] (\xb,-1.0) -- (\xd,-1.0) node[ml]{\code{example.com} ?};
    \node[num] at (-0.5,-1.0) {1};

    \draw[rsp] (\xd,-1.9) -- (\xb,-1.9) node[ml]{\code{93.184.216.34}};
    \node[num] at (-0.5,-1.9) {2};

    \draw[msg] (\xb,-2.9) -- (\xs,-2.9) node[ml]{apertura della connessione TCP};
    \node[num] at (-0.5,-2.9) {3};

    \draw[msg] (\xb,-3.8) -- (\xs,-3.8) node[ml]{\code{GET /dir/file.txt HTTP/1.1}};
    \node[num] at (-0.5,-3.8) {4};

    \draw[rsp] (\xs,-4.7) -- (\xb,-4.7) node[ml]{\code{200 OK} + contenuto del file};
    \node[num] at (-0.5,-4.7) {5};

    \node[font=\Lato\scriptsize\itshape, text=neutral!30!black, anchor=north,
          align=center] at (\xb,-5.8) {il browser interpreta\\ e mostra il contenuto};
    \node[num] at (-1.05,-6.05) {6};
  \end{tikzpicture}
  \caption{Risoluzione DNS e ciclo richiesta--risposta all'apertura di una pagina.}
  \label{fig:browsing}
\end{figure}

Nel dettaglio:

\begin{enumerate}
  \item Il browser estrae dalla URL il \textbf{nome di dominio}
        (\code{example.com}) e chiede a un \term{server DNS} quale sia
        l'indirizzo IP corrispondente.
  \item Il DNS \textbf{risponde} con l'indirizzo IP, ad esempio
        \code{93.184.216.34}.
  \item Il browser \textbf{apre una connessione TCP} verso quell'indirizzo IP,
        sulla porta indicata nella URL (o sulla porta standard, 80 per HTTP e
        443 per HTTPS).
  \item Il browser \textbf{invia la richiesta HTTP}
        (\code{GET /dir/file.txt HTTP/1.1}, con l'header \code{Host}).
  \item Il server elabora la richiesta e \textbf{invia la risposta}, con codice
        di stato, header e corpo.
  \item Il browser \textbf{interpreta} il contenuto ricevuto e lo mostra
        all'utente. Se il contenuto è HTML e fa riferimento ad altre risorse
        (immagini, fogli di stile, script), il browser esegue \textbf{ulteriori
        richieste} per ciascuna di esse.
\end{enumerate}

\info{Il DNS è una rubrica, non un centralino}{
  Il DNS interviene una sola volta, per tradurre il nome in indirizzo IP: il
  traffico HTTP successivo non lo attraversa. Il risultato viene inoltre messo
  in cache dal sistema operativo e dal browser, per cui le visite successive
  allo stesso dominio saltano del tutto questo passaggio.}

\section{HTML: un primo sguardo}

\term{HTML} (\textit{HyperText Markup Language}) è lo standard con cui si
rappresentano i documenti ipertestuali nel Web.

\begin{itemize}
  \item Le pagine web con cui interagiamo nel browser sono documenti definiti
        tramite HTML.
  \item HTML permette di definire pagine con titoli, paragrafi, immagini,
        elenchi, tabelle, form e molto altro.
\end{itemize}

Se HTTP è il \textit{come} si trasferiscono le risorse, HTML è il \textit{che
cosa}: il formato in cui il contenuto è scritto. È l'argomento del prossimo
capitolo.

\section{Riferimenti}

\begin{itemize}
  \item \textbf{Introduction to Web Applications Development}, Carles Mateu.
        Liberamente disponibile su \textit{archive.org} sotto licenza GNU FDL.
        Parti rilevanti: Modulo 1 (\textit{Introduction to Web Applications}).
  \item \textbf{How the web works} --- MDN web docs. \\
        \urlx{https://developer.mozilla.org/en-US/docs/Learn/Getting_started_with_the_web/How_the_Web_works}
  \item \textbf{What are hyperlinks?} --- MDN web docs. \\
        \urlx{https://developer.mozilla.org/en-US/docs/Learn/Common_questions/Web_mechanics/What_are_hyperlinks}
  \item \textbf{What is a URL?} --- MDN web docs. \\
        \urlx{https://developer.mozilla.org/en-US/docs/Learn/Common_questions/Web_mechanics/What_is_a_URL}
  \item \textbf{An overview of HTTP} --- MDN web docs. \\
        \urlx{https://developer.mozilla.org/en-US/docs/Web/HTTP/Overview}
  \item \textbf{What is the difference between webpage, website, web server and
        search engine?} --- MDN web docs. \\
        \urlx{https://developer.mozilla.org/en-US/docs/Learn/Common_questions/Web_mechanics/Pages_sites_servers_and_search_engines}
\end{itemize}
